\chapter{The replica overlap and localization} \label{ch_overlap}

The utlimate goal of this section is to show that the thermodynamic phase transition of section \ref{sec:thermo} is a  localization transition for the polymer. 


To do that, we need fine tools from stochastic analysis. The starting point is Doob-Meyer decomposition, a natural and strong tool to study  stochastic processes in which the process is written as  the sum of a local martingale (the impredictable part) and a bounded variation predictable process (the tamed part).
%The Doob-Meyer decomposition is a natural and strong tool to study  stochastic processes, decomposing it as the sum of a local martingale (the impredictable part) and 
%bounded variation predictable process (a somewhat smooth part), is the starting point.
We start by recalling some martingale properties of the Poisson environment that will prove usefull throughout the following chapters.

%\bfc Trop de theoremes. Simplifier !!! ?? \ec
%%%
\section{The compensated Poisson measure and some associated martingales}
Given the Poisson point process $\om$, we introduce the compensated measure $\bar \om$,
\begin{equation}
\bar \om(\rmd s \rmd x)= \om(\rmd s \rmd x)-\nu\hspace{0.3mm} \rmd s \rmd x,
\end{equation}
and we abreviate its restriction to $(0,t]\times \rd$ by $\bar \om_t$. By definition,  for all function $f(s,x,\om)$ that verifies
\begin{equation} \label{eq:integrabiltyPoissonIntegrand}
\int_{[0,t]\times \mathbb{R}^d} \IP[|f(s,x,\cdot)|] \rmd s \rmd x < \infty,
\end{equation}
the compensated integral of $f$ is given by
\begin{equation}
\int f(s,x,\om) \bar \om_t(\rmd s \rmd x) = \int f(s,x,\om) \om_t(\rmd s \rmd x) - \int_{[0,t]\times\mathbb{R}^d} f(s,x,\om) \nu \hspace{0.3mm} \rmd s \rmd x.
\end{equation}

Furthermore, we say that a function $f(t,x,\om)$ is \emph{predictable}, if it belongs to the sigma-field generated by all the functions $g(t,x,\om)$ that satisfy the following properties:
\begin{enumerate}[label=(\roman*)]
\item for all $t>0$, $(x,\om)\to g(t,x,\om)$ is $\mathcal{B}(\mathbb{R}^d) \times \mathcal{G}_t$-measurable;
\item for all $(x,\om)$, $t\to g(t,x,\om)$ is left continuous.
\end{enumerate}
Then, if the function $f(t,x,\om)$ is predictable, provided that \eqref{eq:integrabiltyPoissonIntegrand} holds and that
\[\int_{[0,t]\times\mathbb{R}^d} \IP[f(s,x,\cdot)^2]\rmd s \rmd x < \infty,\]
the process $t\to \int f \hspace{0.3mm} \rmd \bar \om_t$ is a square-integrable martingale associated to $(\mathcal{G}_t)_{t\geq 0}$, of previsible bracket \cite[Section II.3.]{IkedaWatanabe}
\begin{equation} \label{eq:bracket}
\left\langle \int f \rmd \bar \om \right\rangle_t = \int_{[0,t]\times \mathbb{R}^d} f^2\, \nu \hspace{0.3mm} \rmd s \rmd x.
\end{equation}
The previsible bracket has the property that $ (\int f \hspace{0.3mm} \rmd \bar \om_t)^2 - \langle \int f \hspace{0.3mm} \rmd \bar \om \rangle_t$ is a martingale. In particular,
\begin{equation} \label{eq:MartingaleVarianceFormula}
\mathrm{Var} \left(\int f(s,x,\cdot) \om_t(\rmd s \rmd x)\right) = \int_{[0,t]\times \mathbb{R}^d} \IP[ f(s,x,\cdot)^2] \nu \hspace{0.3mm} \rmd s \rmd x.
\end{equation}
%[Def  $\bar \om$ ??? aussi formule integration de chi s x = rd]


\section{The Doob-Meyer decomposition of $\ln \, Z_t$}
Since $W_t$ is a martingale and $\ln$ is concave, $-\ln(W_t)$ is a submartingale, for which we  
want to get a Doob-Meyer decomposition \cite[Ch.VI]{RogersWilliams}.

In the following, we will use the notation $\Delta_s X := X_s - X_{s-}$ for any c\`adl\`ag process $X$. Let $\zeta_t := e^{\b \om(V_t)}$. As $\om(V_t)$ can be expressed as a sum over $\om_t$, we get by telescopic sum that $\zeta_t = 1 + \int \om_t(\rmd s\rmd x)\Delta_{s} \zeta$, an integral over $\R_+\times\rd$. Averaging over the Brownian path, we also get that $Z_t$ can be expressed as a sum over the process. Therefore, by telescopic sum,
\[\ln Z_{t} = \int \om_t(\rmd s \rmd x) \Delta_s \ln Z.\]
Now, let $(t,x)\in  \R_+ \times\mathbb{R}^d$ be any point of the point process $\om$. As $(t,x)$ is almost surely the only point of $\om$ at time $t$, we can write that
\begin{align*}
Z_{t} &= P\left[ e^{\b}e^{\b \om(V_{t-})} \chi_{t,x}(B) \right] + P\left[ e^{\b \om(V_{t-})} \left(1-\chi_{t,x}(B)\right) \right]\\
& = P\left[ (e^{\b} -1) e^{\b \om(V_{t-})} \chi_{t,x}(B) \right] + P\left[ e^{\b \om(V_{t-})} \right] \\
& =  \left(\lambda(\b)  P_{t-}^{\b,\om} [\chi_{t,x}] + 1\right) Z_{t-}\;.
\end{align*}
Hence,
\begin{equation} \label{eq:logZtIntRepresentation}
\ln Z_{t} = \int \om_t(\rmd s \rmd x) \ln \big(1+\lambda P_{s-}^{\b,\om} [\chi_{s,x}]\big).
\end{equation}
Let $g$ be the function $g(u) = u-\ln(1+u)$, which is positive on $(-1,\infty)$. Then, recalling that $\int_{\rd} \chi_{s,x} dx=r^d$, we obtain  Doob's decomposition 
\begin{equation}\label{eq:DoobDecompLogZ}
-\ln W_t = \lambda \nu r^d t - \ln Z_t = M_t + A_t,
\end{equation}
where the martingale $M$ and the increasing process $A$ are given by
\begin{align} 
M_t & = -\int \bar{\om}_t (\rmd s \rmd x) \ln \big(1+\lambda P_{s-}^{\b,\om} [\chi_{s,x}]\big),\\
A_t & = \lambda \nu r^d t - \int_{[0,t]\times \mathbb{R}^d}  \ln \big(1+\lambda P_{s-}^{\b,\om} [\chi_{s,x}]\big) \, \nu \hspace{0.5mm} \rmd s \rmd x \nonumber\\
& = \int_{[0,t]\times \mathbb{R}^d}  g\big(\lambda P_{s-}^{\b,\om} [\chi_{s,x}]\big) \, \nu \hspace{0.5mm} \rmd s \rmd x. \label{eq:2773}
\end{align}
Moreover, the martingale $M$ is square-integrable, and its bracket is given by
\begin{equation}
\nn
\langle M \rangle_t = \int_{[0,t]\times \mathbb{R}^d} \big[ \ln \big(1+\lambda P_{s-}^{\b,\om} [\chi_{s,x}]\big)\big]^2 \, \nu \hspace{0.5mm} \rmd s \rmd x.
\end{equation}




\section{The replica overlap and quenched overlaps}
%%
Things being clear from the context, we will use the same notation $|A|$ to denote the Lebesgue measure of Borel subset $A$ of $\R^d$ or $\R_+ \times \R^d$.
\begin{definition}
For any two paths $B$ and $\tilde{B}$, we define the \textbf{replica overlap} $R_t(B,\tilde{B})$ as the mean volume overlap of the two tubes around $B$ and $\tilde{B}$ in time $[0,t]$:
\begin{equation}\label{eq:ov1}
R_t(B,\tilde{B})=\frac{1}{tr^d} | V_t(B)\cap V_t(\tilde{B})|.
\end{equation}
In a similar way, we define the \textbf{quenched overlaps} $I_t$ and $J_t$ %[rajouter x ???] 
as\footnote{The product measure ${\gibbs}^{ \otimes 2} = \gibbs \otimes \gibbs$ makes the 2 replicas $B, \tilde B$ independent polymer paths sharing the same environment $\omega$.}
\begin{align} \label{eq:ov2}
I_t &= \frac {1}{r^{d}} %\gibbs \otimes \gibbs \big[ |U(B_t) \cap U(\tilde{B}_t)| \big],\\
{\gibbs}^{ \otimes 2} \big[ |U(B_t) \cap U(\tilde{B}_t)| \big]
,\\ \label{eq:ov3}
J_t &=   \frac {1}{tr^{d}}  % \gibbs \otimes \gibbs \big[ R_t(B,\tilde{B})\big] = \frac{1}{tr^d} \int_0^t \gibbs \otimes \gibbs |U(B_s) \cap U(\tilde{B}_s)| \, \rmd s.
{\gibbs}^{ \otimes 2} \big[ |V_t(B) \cap V_t(\tilde{B})| \big],
\end{align}
\end{definition}

The  variable $I_t$ stands for the expected volume of overlap around the endpoints of two  independent polymer paths, while $J_t$ is the expected volume of overlap during the time interval $[0,t]$.
Note that both $\int_0^t I_s \rmd s$ and $t J_t$ represent an expected volume of overlap in time $[0,t]$, but they will emerge from different circumstances. 
Similar to  $\chi_{s,x}= {\mathbf 1} \{x \in U(B_s) \}$ from definition \eqref{def:chitx}, we write for short 
\begin{equation}\nn
\tilde \chi_{s,x}= {\mathbf 1} \{x \in U(\tilde B_s) \} .
\end{equation}
Writing 
$$
|U(B_t) \cap U(\tilde{B}_t)|  = \int_{\rd}  \chi_{t,x}\tilde \chi_{t,x} dx,
$$
we derive two useful formulas:
\begin{equation}\label{eq:useful}
I_t=\frac{1}{r^{d}}  \int_{\rd} \gibbs (\chi_{t,x})^2 dx\;,\qquad J_t = \frac{1}{r^{d}} \int_{\rd} dx \frac 1t  \int_0^t \gibbs (\chi_{s,x})^2 ds\;.
\end{equation}

\medskip

For better comparisons, we have normalized all quantities in \eqref{eq:ov1}, \eqref{eq:ov2} and \eqref{eq:ov3}  in such a way that
$$
0 \leq R_t, I_t, J_t \leq 1 \;,
$$
so that we can -- and we will -- view each of them as a \emph{localization index}:
\begin{itemize}
\item   $R_t$ close to 1 means that the two fixed paths $B, \tilde B$ are close on the interval $[0,t]$;
\item $I_t$  close to 1 means that the endpoints of two independent samples of the polymer measure are typically close one from the other;
\item $J_t$ close to 1 means that the paths of two independent polymers are close all along the time interval.
\end{itemize}
The second case corresponds to \emph{endpoint localization} whereas the third one is  \emph{path localization}. 
Mathematically, the quantity $I_t$ appears via It\^o's calculus (stochastic differentiation) and $J_t$ via Malliavin calculus (integration by parts).
On the contrary, small values of these indices correspond to
absence of localization: it means that the polymer spreads more or less uniformly in space without particular preference.


%
%\begin{remark} 
%These two quantities can be seen as indicators of [the impact of the environment over] [the localization properties] the polymer. Indeed, for high values of $|\b|$, two 
% independent polymers can be tempted to go to the same parts of the environment that feature high rewards, so that they could be encouraged to stay close 
%to one another for non-negligeable [linear ???] time, this leading to $\int_0^t I_s \rmd s$ and $J_t$ to take higher values that in the diffusive case.
%\end{remark} 

\begin{remark} 
When $\b = 0$, the Gibbs measure $\gibbs$ reduces to Wiener measure $P$, so that $J_t = \int_0^t I_s \rmd s = t P^{\otimes 2} [R_t(B,\tilde{B})]$, and
\[P^{\otimes 2} [R_t(B,\tilde{B})] = \frac{1}{t} \int_0^t P^{\otimes 2} \big[|U(B_s)\cap U(\tilde{B}_s)|\big] \rmd s \leq \frac{r^d}{t} \int_0^t P\left(B_{2s}\in U(0)\right) \rmd s \underset{t\to 0}{\longrightarrow} 0 .\] Thus, \[\lim_{t\to\infty}  J_t = \lim_{t\to\infty} \frac{1}{t} \int_0^t I_s \rmd s = 0,\] 
This indicates that, in absence of interaction with the inhomogeneous medium,  there is no localization of the Brownian path.
\end{remark} 
%%%
We now come to the core of the section: localization results.
%%%
\section{Endpoint localization}
%
A naive prediction is  that, for small $\b$,  the polymer is a small perturbation of Brownian motion for small $\b$, with a comparable behavior, whereas  for large $\b$ localization takes place and the limits are nonzero.
A first theorem shows that this is the case for the quantity $I_t$:
\begin{theorem} \label{th:loca} The following equivalence holds for $\b \neq 0$:
\begin{equation} \label{eq:intIsEqSD}
W_\infty = 0 \iff \int_0^\infty I_s \, \rmd s = \infty, \quad \IP\text{-a.s.}
\end{equation}
In particular, the above integral is a.s. finite for $\b \in (\bar{\b}_c^-, \bar{\b}_c^+)$, and a.s. infinite for $\b < \bar{\b}_c^-$ or $\b > \bar{\b}_c^+$.
Moreover, we have: 
\begin{align} 
&\lim_{t\to\infty} \frac{1}{t} \int_0^t I_s \, \rmd s = 0 \  \textup{ if } \b \in [\b_c^-,\b_c^+]\cap \mathbb{R},\label{eq:AsymptIntIs1}\\
&\liminf_{t\to\infty} \frac{1}{t} \int_0^t I_s \, \rmd s > 0 \  \textup{ if } \b \in \mathbb{R}\setminus [\b_c^-,\b_c^+]. \label{eq:AsymptIntIs2}
\end{align}
\end{theorem}
\begin{remark} Note that  $\int_0^\infty I_s \rmd s$ is a.s. finite or a.s. infinite,  by the dichotomy between strong and weak disorder. Similarly, 
%the asymptotic behavior of its mean is linked to the free energy bound [??? how to call it.Very strong vs not so strong disorder?]. In fact, we will prove
strict positivity of $\,\liminf_{t\to\infty} t^{-1} \int_0^t I_s ds$  is equivalent to low temperature region $\psi(\b, \nu)>0$ in \eqref{eq:highlow}.
In fact, we will prove
in \eqref{eq:encadrementAtIs}--\eqref{eq:freeEnergyOverAt}
 that under strong disorder, there exist $c_1,c_2 \in (0,\infty)$, such that
\begin{equation} \label{eq:EncadrementFreeEnergy}
c_1 \int_0^t I_s \, \rmd s \leq  - \ln W_t \leq c_2 \int_0^t I_s \, \rmd s, \quad \text{for large t}, \quad \IP{-a.s.}
\end{equation}
\end{remark}
\begin{proof} (Theorem \ref{th:loca})
First observe that one can easily derive \eqref{eq:AsymptIntIs1} and \eqref{eq:AsymptIntIs2} from \eqref{eq:intIsEqSD} and \eqref{eq:EncadrementFreeEnergy}.

To show \eqref{eq:EncadrementFreeEnergy}, we will relate $\int_0^t I_s\rmd s$ to the variables $M_t$ and $A_t$ of the Doob decomposition \eqref{eq:DoobDecompLogZ} as follows. From %the left-hand side of 
\eqref{eq:useful}, we have
\begin{equation} \label{eq:IntIsEqGibbsSquared} \int_{[0,t]\times\mathbb{R}^d} \left(P_s^{\b,\om}[\chi_{s,x}]\right)^2 \rmd s \rmd x = 
%\int_0^t P_s^{\b,\om}\otimes P_s^{\b,\om} [U(B_s)\cap U(\tilde{B}_s)] \rmd s  =
 r^d \int_0^t I_s \, \rmd s .
\end{equation}
Then, looking at the behavior of $g(u)$ in \eqref{eq:2773} 
and $\ln(1+u)$ around $0$, it is clear that there are two constants $c_1,c_2>0$, depending only on $\b$ and $\nu$, such that
\[ c_1 u^2 \leq \nu g(u) \leq c_2 u^2, \quad \quad \nu\ln(1+u)^2 \leq c_2 u^2,\]
for all $u$ in $[0,\lambda]$ when $\b\geq 0$, (resp. $[\lambda,0]$ when $\b\leq 0$). Together with \eqref{eq:IntIsEqGibbsSquared}, this implies that %[??? crochet defini?]
\begin{align}\label{eq:encadrementAtIs}
c_1 \int_0^t I_s \rmd s \leq & A_t \leq  \, c_2 \int_0^t I_s \rmd s,\\
\langle & M \rangle_t \leq  \, c_2 \int_0^t I_s \rmd s.\label{eq:majCrochetMtIs}
\end{align}

We then recall two results about martingales. (These facts for the discrete martingales are standard (e.g. \cite[p. 255, (4.9),(4.10)]{Dur95}. It
is not difficult to adapt the proof for the discrete setting to our case.)
Let $\varepsilon >0$, then
\begin{align} &\langle M \rangle_\infty < \infty \implies (M_t)_{t\geq 0} \text{ converges\; a.s.}, \label{eq:BracketImpliesCV}\\  
&\langle M \rangle_\infty = \infty \implies \lim_{t\to\infty} M_t / \langle M \rangle_t^{\frac{1+\varepsilon}{2}} = 0 \quad {\rm a.s.}. \label{eq:InfiniteBracket}
\end{align}
Consequently, we get that $\IP$-almost surely,
\begin{align} \nn
\int_0^\infty I_s \rmd s < \infty & \iff A_\infty < \infty,  \langle M \rangle_\infty < \infty\\ \nn
& \implies  A_\infty < \infty, \, \lim_{t\to\infty} M_t \text{ exists and is finite}\\ \nn
& \implies W_\infty > 0,
\end{align}
where the last implication comes from the Doob decomposition \eqref{eq:DoobDecompLogZ}. By contraposition, this proves the first implication of \eqref{eq:intIsEqSD}.
The reverse implication will follow from the arguments below.

The next step is to show \eqref{eq:EncadrementFreeEnergy}, so we now suppose that we are in the strong disorder setting.
Using \eqref{eq:encadrementAtIs}, we see that it is enough to show that
\begin{equation} \label{eq:freeEnergyOverAt}
\lim_{t\to\infty} \frac{- \ln W_t}{A_t} = 1, \quad \IP\text{-a.s.}
\end{equation}
or equivalently by the Doob decomposition, that 
\begin{equation}\label{eq:MtOnAtCV0}
\lim_{t \to \infty} \frac{M_t}{A_t} = 0.
\end{equation}

As we just proved, the strong disorder implies that $\int_0^\infty I_s \rmd s = \infty$, which in turn implies that  $A_\infty = \infty$ by \eqref{eq:encadrementAtIs}. Thus, in the case that $\langle M \rangle_\infty < \infty$, the martingale $M_t$ converges so that the condition \eqref{eq:MtOnAtCV0} directly holds. When $\langle M \rangle_\infty = \infty$, this is still true as
\[\frac{M_t}{A_t}=\frac{M_t}{\langle M \rangle_t} \frac{\langle M \rangle_t}{A_t} \underset{t\to\infty}{\longrightarrow} 0,\]
by (\ref{eq:encadrementAtIs}, \ref{eq:majCrochetMtIs},  \ref{eq:InfiniteBracket}).

Finally, what is left to demonstrate is that $\IP$-almost surely   $W_\infty = 0$ 
on the event $\int_0^\infty I_s \rmd s=\infty$. In fact, we showed that this event implies both the limit \eqref{eq:freeEnergyOverAt} and $A_\infty = \infty$, so that $W_t \to 0$.
%$\IP$-almost surely,
% if $\int_0^\infty I_s \rmd s=\infty$ then $W_\infty = 0$. In fact, we showed that this first condition implies both the limit \eqref{eq:freeEnergyOverAt} 
%and $A_\infty = \infty$, so that $W_t \to 0$.
\end{proof}
%%%%%%%%%%%%%

{\bf Endpoint localization}: 
As in the discrete case, we can interpret the results from the present subsection, in terms
of localization for the path. Indeed, it is proven in Sect. 8 of \cite{CY05}, that for some constant $c_1$, 
\begin{equation} \label{eq:overlap-favorite}
c_1 \sup_{y \in \rd} \gibbs [ B_s \in U(y)]^2 \leq { \gibbs}^{\otimes 2} \left[ \big \vert U(B_s) \bigcap U(\tilde B_s)\big \vert \right]
\leq \sup_{y \in \rd} \gibbs [ B_s \in U(y)]
\end{equation}
% [C'EST \eqref{eq:boundOnRt} ????]
(in fact, the inequality on the right is trivial, and the one on the left is the combination of \eqref{eq:num-1} and \eqref{eq:num-2}).


The maximum appearing in the above bounds should be viewed as the probability of the favorite location for $B_s$, under the polymer measure $\gibbs$; for $s=t$, the supremum is called the probability of the \emph{favorite endpoint}, and the maximizing $y$ is
the \emph{location of the favorite endpoint}.
Both Theorem \ref{th:loca} and Theorem \ref{th:asymptOfJt} are precise statements that the polymer localizes in the strong disorder
regime in a few specific corridors of width ${\cal O}(1)$, but spreads out in a diffuse way in the
weak disorder regime. If $\psi(\b,\nu)>0$, the Cesaro-limit of  probability of the {favorite endpoint} is strictly positive.


\section{Favorite path and path localization}



%
%[CHANGE notation $\left(\frac{\partial \psi}{\partial \b}\right)_+  $ and DEFINE it ????]
%
Recall that the excess free energy $\psi$ from 
\eqref{def:excessfreeenergy} is the difference of a smooth function and a convex function. Hence its right-derivative, resp. left-derivative, 
\begin{equation}\nn
 \left( \frac{\partial \psi}{\partial \b}\right)_+ (\b,\nu) = \lim_{\b' \searrow \b} \frac{\psi(\b',\nu)-\psi(\b,\nu)}{\b'-\b}, \qquad
 \left(   \frac{\partial \psi}{\partial \b}\right)_- (\b,\nu) = \lim_{\b' \nearrow \b} \frac{\psi(\b',\nu)-\psi(\b,\nu)}{\b'-\b},
\end{equation}
exists for all $\b$ and all $\nu$, and satisfy 
$ \left( \frac{\partial \psi}{\partial \b}\right)_+ (\b,\nu) \leq  \left(\frac{\partial \psi}{\partial \b}\right)_- (\b,\nu)$. For the same reason as above, $\psi(\cdot,\nu)$ is differentiable except on a set
which is at most countable, and we can write
 \begin{equation}\nn
 \left(  \frac{\partial \psi}{\partial \b}\right)_+ (\b,\nu) = \lim_{\b'\geq \b } 
  \frac{\partial \psi}{\partial \b} (\b',\nu) ,
\end{equation}
where the limit is over differentiability points $\b'$ tending to $\b$ by larger values. A similar statement holds for the left-derivative.
For further use, we note that for all fixed $\nu$, $\psi(\cdot, \nu)$ is absolutely continuous again for the same reason as above.

Now, we turn to the properties of the replica overlap $J_t$. The key fact is the following proposition:
\begin{proposition} \label{prop:boundJt}
There exist two constants $c_1,c_2 \in (0,\infty)$, depending only on $\b$ and $\nu$, such that
\begin{align} \label{eq:bound_PJt1}
\forall \b>0, & \quad c_1 \,  \left(  \frac{\partial \psi}{\partial \b}\right)_+\leq \liminf_{t\to\infty} \,  \IP[J_t] \leq \limsup_{t\to\infty} \,  \IP[J_t] \leq c_2\,\left(  \frac{\partial \psi}{\partial \b}\right)_-, \\
\label{eq:bound_PJt2} \forall \b<0, & \quad -c_1\,  \left(  \frac{\partial \psi}{\partial \b}\right)_-\leq \liminf_{t\to\infty} \,  \IP[J_t] \leq \limsup_{t\to\infty} \, \IP[J_t] \leq -c_2 \,
 \left(  \frac{\partial \psi}{\partial \b}\right)_+. 
\end{align}
\end{proposition}
\begin{proof} It is not difficult to see from the definition that $t J_t = \iint_{[0,t]\times \mathbb{R}^d} \gibbs[\chi_{s,x}]^2 \rmd s \rmd x$. Hence, using equation \eqref{eq:psi'} and the fact that
$e^{-|\b|} \leq 1+\lambda \gibbs[\chi_{s,x}] \leq e^{|\b|}$, we obtain that
\begin{equation} \label{eq:bound_dpsit}
\lambda \nu e^{\b-|\b|} t \IP[J_t] \leq t \frac{\partial}{\partial \b}\psi_t(\b,\nu) \leq \nu \lambda e^{\b + |\b|} t \IP[J_t].
\end{equation}
Moreover, the excess free energy writes $\psi(\b,\nu)=\nu \lambda(\b)r^d- p(\b, \nu)$, where $p(\b,\nu)$ is a convex function, defined as the limit, for $t\to\infty$, of the convex functions
$p_t(\b,\nu) = \frac{1}{t}\IP[\ln Z_t].$
By convexity properties, we know that
\[\left(  \frac{\partial p}{\partial \b}\right)_-\leq \liminf_{t\to \infty} \frac{\partial p_t}{\partial \b } \leq \limsup_{t\to \infty} \frac{\partial p_t}{\partial \b } \leq
\left(  \frac{\partial p}{\partial \b}\right)_+,
\]
which in turns implies that
\begin{equation} \left(  \frac{\partial \psi}{\partial \b}\right)_+\leq \liminf_{t\to \infty} \frac{\partial \psi_t}{\partial \b } \leq \limsup_{t\to \infty} \frac{\partial \psi_t}{\partial \b } \leq
\left(  \frac{\partial \psi}{\partial \b}\right)_-
.\end{equation}
The proposition is then a consequence of \eqref{eq:bound_dpsit} and these last inequalities.
\end{proof}
With this proposition, we can give a characterization of the critical values $\b_c^\pm$, in terms of the asymptotics of the overlap:
\begin{theorem} \label{th:asymptOfJt}
For all $\b \in [\b_c^-,\b_c^+]\cap \mathbb{R}$,
\begin{equation}
\lim_{t\to\infty}  \IP[J_t] = %\psi(\b) = 
0.
\end{equation}
Furthermore,
\begin{align} 
\b_c^+ = \sup \{\b' \geq 0 : \forall \b \in [0,\b'], \, \lim_{t\to\infty} \IP[J_t] = 0\}= \inf\{ \b>0: \liminf_{t\to\infty} \IP[J_t] > 0\},\\
\b_c^- = \inf \{\b' \leq 0 : \forall \b \in [\b',0], \, \lim_{t\to\infty}  \IP[J_t] = 0\}= \sup\{ \b<0: \liminf_{t\to\infty} \IP[J_t] > 0\}.
\end{align}
\end{theorem}
\begin{proof} We will focus on the $\b\geq 0$ case, but the same arguments can be applied to the $\b \leq 0$ case. Define
\[ \delta_c^+ = \sup \{\b' \geq 0 : \forall \b \in [0,\b'], \, \lim_{t\to\infty}  \IP[J_t] = 0\},\]
so what we need to show in particular is $\b_c^+ = \delta_c^+$.

To prove the first claim of the theorem, from which $\b_c^+ \leq \delta_c^+$ follows directly, it is enough, using \eqref{eq:bound_PJt1}, to verify that
\begin{equation}
\forall \b \leq \b_c^+, \quad \left( \frac{\partial \psi}{\partial \b}\right)_-(\b,\nu) = 0.
\end{equation}
This property is true when $\b\in [0,\b_c^+)$, as $\psi$ is constant and set to $0$ in this interval. To prove that it extends to $\b_c^+$ if $\b_c^+ < \infty$, observe that $\psi$ is minimal at $\b_c^+$, so that
\[ \left( \frac{\partial \psi}{\partial \b}\right)_- (\b_c^+,\nu) \leq 0 \leq  \left(\frac{\partial \psi}{\partial \b}\right)_+ (\b_c^+,\nu).\]
As we saw earlier that $\frac{\partial \psi}{\partial \b_+} \leq  \frac{\partial \psi}{\partial \b_-}$ always holds, we finally get that $\frac{\partial \psi}{\partial \b_- }(\b_c^+) = \frac{\partial \psi}{\partial \b_+ }(\b_c^+)= 0$.

%[??? (i) Did we mention that $\b_c^+ < \infty$ ? As it is not true for $b_c^-$, we need to keep the condition here in the proof (ii) Tell what happens for the
% derivatives at the critical point before here?]


We now prove $\b_c^+ \geq \delta_c^+$. Let $\b > \b_c^+$ be finite,
%[$\b_c <\infty$ ???] ???, PAS toujours pour $\bar{\b}_-$.
so that, by definition, $\psi(\b,\nu)>0$. As $\psi$ is absolutely continuous 
%a continuous function of $\b$ [sufficient ???] and since 
with $\psi(0,\nu) = 0$,  one can write
\[\psi(\b,\nu) = \int_0^\b \frac{\partial \psi}{\partial \b_+ }(\b',\nu)\rmd \b' > 0,\]
which implies that there exists some $\b'\leq \b$ such that $\frac{\partial \psi}{\partial \b+ }(\b',\nu) > 0$. By equation \eqref{eq:bound_PJt1}, we get that $\liminf_{t\to\infty} \frac{1}{t} \IP[J_t(\b)] > 0$, hence $\b \geq \delta_c^+$. As it is true for all $\b > \beta_c^+$, we obtain that $\b_c^+ \geq \delta_c^+$.
\end{proof}

\medskip
\paragraph{The favorite path.}
Let $\mathcal{M}$ be the set of integer-valued Radon measures on $\mathbb{R}_+\times \mathbb{R}^d$, equipped with the sigma-field $\mathcal G$ generated by the variables $\om(A)$, $A\in \mathbb{R}_+\times \mathbb{R}^d$, so that we will consider $\om$ as a process of the probability space $(\mathcal{M},\mathcal{G},\IP)$.
It is possible to define, for all fixed time horizon $t>0$, a measurable function
\begin{equation} \begin{array}{cccc}
\Yt:&[0,t]\times\mathcal{M} & \to & \mathbb{R}^d\\
&(s,\om) & \mapsto & \Yt_s,
\end{array}
\end{equation}
which satisfies the property that, $\IP$-almost surely,
\begin{equation}
\forall s\in[0,t],\quad \gibbs \left(B_s \in U\left(\Yt_s\right)\right)=\max_{x\in\mathbb{R}^d} \gibbs \big(B_s \in U(x) \big).
\end{equation}
The reader may refer to \cite{CYBMPO2} for a proper (and rather technical) definition of $\Yt$. 

Here, the path $s\to \Yt_s$ stands for the "optimal path" or the "favorite path" of the polymer, although this path is neither necessarily continuous, nor necessarily unique.
Similarly to what we have done previously, we define the overlap with the favorite path $R_t^*$ as the fraction of time any path $B$ stays next to the favorite path:
\begin{equation}
R_t^* = R_t^*(B,\om) = \frac{1}{t} \int_0^t \mathbf{1}_{B_s\in \, U(\Yt_s)} \rmd s.
\end{equation}

As discussed before, a question of interest is the asymptotic behavior, as $t\to\infty$, of $R_t$ and of $R_t^*$, and we will see in Theorem \ref{th:localization} that they are related. In particular, we are interested in determining the regions were one can prove positivity in the limit of these quantities, which can be seen as localization properties of the polymer. 

%Indeed, if the limit inferior of $R_t^*$ were to be positive with some positive probability [est-ce le cas???], it would tell us that there is a specific path, depending only on 
%the environment, to which the polymer is susceptible to stay close for a positive fraction of its time. [??? This is in contrast to Brownian motion, where]. 

Recall the notations
\begin{equation} \nn
{\cal D} = \{ (\b,\nu): \psi(\b, \nu)=0 \}, \quad {\cal L}  = \{ (\b,\nu): \psi(\b, \nu)>0\},
\end{equation}
of the high and low temperature regions, which are delimited by the critical lines $\b_c^-(\nu)$ and $\b_c^+(\nu)$ (cf. Definition \ref{defdef..}).
We saw in theorem \ref{th:asymptOfJt} that in the ${\cal D}$ region, $\lim_{t\to\infty} \IP[J_t] = 0$. On the other hand, Proposition \ref{prop:boundJt} tells us that the limit inferior of $\IP[J_t]$ is always positive in the region $\mathcal{L}'$, where
\begin{equation} \label{eq:Lregion}
\mathcal{L}' = \left\{ \b > 0,\nu > 0 :  \left(  \frac{\partial \psi}{\partial \b}\right)_+ > 0 \right\}\cup\left\{ \b < 0, \nu > 0 :  \left(  \frac{\partial \psi}{\partial \b}\right)_- < 0 \right\}.
\end{equation}
From the preceeding considerations, we know that
\begin{equation}% \mathcal{L}'\cap \mathcal{D} = \emptyset,
\mathcal{L}' \subset \mathcal{L}\;,
\end{equation}
and a still open question is whether 
$\mathcal{L}' = \mathcal{L}$ or not.
%[Say that L would then be con ??? confirm this speech is ok. Introduce L and D before?]

\begin{remark}
It is a direct consequence of the monotonicity of $\psi$ (point \ref{eq:p3} of proposition \ref{prop:p}), that the inequalities on the derivatives of $\psi$ appearing in \eqref{eq:Lregion}, once replaced by large inequalities, are always verified:
$$ \b<0 \implies  \left(  \frac{\partial \psi}{\partial \b}\right)_- \leq 0,
\quad
 \b>0 \implies  \left(  \frac{\partial \psi}{\partial \b}\right)_+ \geq 0.$$
\end{remark}

We first state some results about the localized region:
\begin{proposition}
\begin{enumerate}[label=(\roman*)]
\item For any fixed $\nu>0$ and for large enough positive $\b$, $(\b,\nu) \in \mathcal{L}'$.
\item For all $(\b,\nu) \in \mathcal{L}'$, $\liminf_{t\to\infty} \IP[J_t] > 0$.
%\item ??? points d accroissement [cf. Th. 3.3.1 ?????]
\end{enumerate}
\end{proposition}
\begin{proof}
To prove (i), we use a result of \cite[Th. 2.2.2.(b)]{CY05}
%\cite[Th. 5.3.1]{CY13}: 
where it is shown that there exists a positive constant $C_1=C_1(r,\nu,d)$, such that, for fixed $\nu, r$ and $\b$ large enough,
\[p(\b,\nu) \leq C_1 \lambda^{1/2}.\]
By convexity of $p$ in $\b$, we get that for large enough $\b$,
%\begin{align*}\left(  \frac{\partial p}{\partial \b}\right)_+ &\leq p(\b+1,\nu) - p(\b,\nu)\\
%& \leq C_1 e^{(\b+1)/2}\\
%& < \nu r^d e^\b,
%\end{align*}
$$\left(  \frac{\partial p}{\partial \b}\right)_+ \leq p(\b+1,\nu) - p(\b,\nu) 
 < \nu r^d e^\b,
$$
so that
$\left(  \frac{\partial \psi}{\partial \b}\right)_+$ is indeed positive when $\b$ is big enough.

The property (ii) is given by proposition \ref{prop:boundJt}.
\end{proof}
%%
\begin{remark}
We stress on how strong is the above claim (ii). For $(\b,\nu) \in \mathcal{L}'$, it implies that   there exist $C>0, \delta>0$ such that 
$$\liminf_{t\to\infty} \IP  {\gibbs}^{ \otimes 2} \big[ R_t(B, \tilde{B}) \geq \delta \big] \geq C\;.$$
In contrast, if $\b=0$, there is  some $C'>0$ such that
$$
P^{\otimes 2}  \big[ R_t(B, \tilde{B}) \geq \delta \big] \leq e^{-C't}
$$
for all large enough $t$.
\end{remark}
%%
Observe that the properties of $R_t$ and $R_t^*$ are comparable in the following sense:
\begin{theorem} \label{th:localization}
There exists a constant $c=c(d,r)$ in $(0,1)$, such that
\begin{equation} \label{eq:2to1}
c\left(\IP\gibbs \big[R_t^*\big] \right)^2 \leq \IP[J_t] \leq r^d \, \IP \gibbs\big[R_t^*\big].
\end{equation}
In particular, we get that for all $\b\in\mathcal{L}'$,
\begin{equation} \label{eq:localization}
\liminf_{t\to\infty} \IP \gibbs \big[R_t^*\big] \geq r^{-d}\liminf_{t\to\infty} \IP[J_t] >0.
\end{equation}
\end{theorem}
\begin{remark}
Note that equation \eqref{eq:localization}  gives another feature of path localization of the polymer in the $\mathcal{L}'$ region: we can find a "path" depending only on the environment (here, we found that $\Yt$ does the job) such that the {\rm expected proportion of time the random polymer spends in the neighborhood of that "path" 
is bounded away from 0} as $t \to \8$. Under the Gibbs measure, the random polymer sticks to that particular "path". Even though that "path" is not smooth -- in fact, it has long jumps -- it is an interesting object which sumarizes the attractive effect  of the medium.
\end{remark}
%%%%%
\begin{proof}
To prove the first point, it is enough to show that there is a $c\in(0,1)$, such that
\begin{equation} \label{eq:boundOnRt}
c\gibbs\big[R_t^*\big]^2 \leq J_t \leq r^d\gibbs\big[R_t^*\big],
\end{equation}
the proposition being then a simple consequence of Jensen's inequality. 
For the right-hand side inequality of \eqref{eq:boundOnRt}, observe that by Fubini's theorem,
\begin{align*}
\gibbs\otimes\gibbs\big[R_t\big] &=\frac{1}{t}\int_0^t \int \gibbs[\chi_{s,x}]^2 \rmd s \rmd x.\\
& \leq \frac{1}{t}\int_0^t \max_{x\in\mathbb{R}^d} \gibbs\left( B_s \in U(x)\right)\rmd s \times \int_{\mathbb{R}^d} \gibbs\big[\chi_{s,x}\big] \rmd x\\
& = r^d \gibbs\left[ R_t^* \right].
\end{align*}

In order to obtain the left-hand side inequality, let $r_d$ denote the radius of the ball $U(x)$ and $y$ be any point of $\mathbb{R}^d$. By Cauchy-Schwarz's inequality,
\[\left(\int_{\mathrm{B}(y,r_d/2)}  \gibbs\big(B_s \in U(z) \big) \rmd z \right)^2 \leq \big|\mathrm{B}(y,r_d/2)\big| \int_{\mathbb{R}^d} \gibbs\big(B_s \in U(z)\big)^2,\]
and since for every $z$ in $\mathrm{B}(y,r_d/2)$, the ball $\mathrm{B}(y,r_d/2)$ is included in $U(z)$, this inequality leads to the following:
\begin{align}\nn
\int_{\mathbb{R}^d} \gibbs\big(B_s \in U(z)\big)^2 \rmd z &\geq \frac{2^d}{r^d} \left(\int_{\mathrm{B}(y,r_d/2)}  \gibbs\big(B_s \in \mathrm{B}(y,r_d/2) \big) \rmd z \right)^2\\
&= \frac{r^d}{2^d} \, \gibbs\big(B_s \in \mathrm{B}(y,r_d/2)\big)^2. \label{eq:num-1}
\end{align}

Now, let $c'=c'(d)$ be the minimal number of copies of $\mathrm{B}(y,r_d/2)$ necessary to cover $U(y)$. Then, by additivity of $\gibbs$,
\begin{equation} \label{eq:num-2}
\max_{y\in\mathbb{R}^d} \gibbs\big(B_s\in U(y)\big) \leq c' \max_{y\in\mathbb{R}^d} \gibbs\big(B_s\in \mathrm{B}(y,r_d/2)\big).
\end{equation}
Putting things together and integrating on $[0,t]$, we finally get that
\[ c\,\frac{1}{t} \int_0^t \max_{y\in\mathbb{R}^d} \gibbs\big(B_s\in U(y)\big)^2 \rmd s \leq  \frac{1}{t}\int_0^t \int \gibbs[\chi_{s,x}]^2 \rmd s \rmd x,\]
where $c=(c')^{-2} {r^d}/{2^d}$, from which the left-hand side inequality of \eqref{eq:boundOnRt} can be obtained by applying Jensen's inequality with probability measure $\rmd s /t$ on [0,t].

The second point of the theorem is then a consequence of the first point and proposition \ref{prop:boundJt}.
\end{proof}
\begin{remark}
Formulas like \eqref{eq:2to1} and \eqref{eq:boundOnRt} can be called 2-to-1 formulas since they relate quenched expectations for two independent polymers to
expectations for only one polymer (and involving the optimal path). As we have seen from the computations, stochastic analysis brings in second moments, involving 2 replicas of the polymer path. Then, using such formulas, the information is reduced to one polymer path interacting with the favorite path.
\end{remark}
